Este trabalho apresenta um laboratório experimental de detecção de intrusão em redes com foco na comparação entre Random Forest e Rede Neural Profunda (DNN) como classificadores supervisionados, avaliando os efeitos da qualidade dos dados e da redução de dimensionalidade por \textit{Mean Decrease in Impurity} (MDI) sobre os \textit{datasets} CICIDS2017 e LycoS-IDS2017. A hipótese inicial previa que a DNN, pela sua maior capacidade de representação, apresentaria melhor generalização para ataques ausentes no treinamento; contudo, essa expectativa não se confirmou. Diante dessa limitação, um Autoencoder treinado exclusivamente com tráfego benigno foi investigado como alternativa não supervisionada para a detecção de ataques desconhecidos. O Random Forest superou a DNN em todos os cenários avaliados, atingindo Macro F1 de 0,97 no LycoS-IDS2017 contra 0,87 da DNN, com tempo de inferência até 5,18 vezes menor; a DNN apresentou dificuldades em classes minoritárias e revelou maior sensibilidade à qualidade dos dados, com ganho de 0,20 pontos de Macro F1 ao substituir o CICIDS2017 pelo LycoS-IDS2017, frente a 0,11 do Random Forest. A redução de dimensionalidade por MDI prejudicou todos os modelos por favorecer atributos das classes majoritárias em detrimento de características discriminativas para ataques raros. O Autoencoder alcançou Macro F1 de 0,89 e \textit{Recall} de ataque de 0,82 no LycoS-IDS2017, demonstrando viabilidade parcial para cenários sem rótulos, embora sem superar os classificadores supervisionados em ataques conhecidos. Conclui-se que o Random Forest com dados curados é a solução mais adequada para produção com ataques conhecidos, que a qualidade do \textit{dataset} mostra-se mais determinante do que a complexidade da arquitetura, e que a detecção de anomalias desconhecidas permanece um desafio em aberto.

\noindent\textbf{Palavras-chave:} Detecção de intrusão. \textit{Random Forest}. Redes neurais profundas. \textit{Autoencoder}. CICIDS2017.
